\section{Introduction}\label{sec: intro}\noindent
Spatio-temporal outputs are inherent in geotechnical systems, where stress, pore pressure, and deformation evolve over space and time under changing loading and environmental conditions \parencite{phoon2020,qin2021,shi2021,chen2023,phoon2023,gatto2023}. These responses are essential for assessing deformation, stability, and fluid flow. Numerical methods, such as finite element method (FEM), are commonly used to obtain such outputs based on input soil parameters, construction sequences, and a discretised mesh. Predicted outputs can range from scalar values to full spatial fields. However, high-fidelity FEM are computationally expensive, particularly for real-time or probabilistic applications like uncertainty quantification (UQ) or digital twins \parencite{kapteyn2021,torzoni2024,tian2025}, which often require numerous model evaluations. Running FEM for each likelihood evaluation in processes such as MCMC-based inversion is prohibitively costly. Surrogate models (meta-models) offer a more efficient alternative by approximating system responses at a lower computational cost, enabling faster evaluations without repeatedly running the full FEM. This makes it possible to perform complex analyses, such as UQ and inversion tasks, that would otherwise be too expensive or time-consuming. 

Despite the efficiency of metamodels, many existing machine learning methods treat spatial and temporal outputs separately \parencite{bardet2009,low2014,sotiropoulos2016,zhao2021,phoon2022,chen2024}, often neglecting the inherently coupled, high-dimensional, and dynamic characteristics of geotechnical problems. Machine learning models have recently become popular tools for surrogate modelling in geotechnical engineering \parencite{zhang2022,wang2023,wang2025-1,wang2025-2, tian2025,zhang2025}. Methods such as sequential layered neural networks \parencite{chen2024,yaghoubi2024}, support vector machines \parencite{goh2007,tinoco2014}, and polynomial chaos expansions \parencite{pan2017,pan2020,mohammadi2023} typically address either spatial or temporal variations in isolation. In our earlier work \parencite{yang2024}, a two-step reduced-order modelling (ROM) framework was developed to capture correlations among geotechnical outputs. In this framework, principal component analysis (PCA) is employed to extract the dominant physical modes embedded within the outputs. As PCA is inherently a matrix-based technique, requiring outputs to be represented in vector form. As a result, the constructed models are restricted to a single loading stage and a single type of output, such as the excavation wall or pile deflections at a fixed stage. However, many geotechnical problems involve outputs generated over multiple stages, exhibiting strong temporal dependencies and encompassing diverse quantities of interest (e.g. stresses or displacements). These outputs may take different forms, including both field and vector data, and evolve with time, which further complicates the direct application of PCA-based dimensionality reduction. 

An effective surrogate for geotechnical outputs should therefore integrate both spatial characteristics and temporal dependencies. Several machine learning approaches have been explored for spatial and temporal modelling. Reduced-order models and convolutional neural networks (CNNs) are widely employed to represent spatial data \parencite{liu2019,smith2019,tian2020,recanatesi2021,lui2021,huang2022,xu2023,raju2024}, while temporal dependencies are commonly addressed using recurrent architectures such as long short-term memory (LSTM) networks \parencite{zhou2022,wang2022,jiang2023,de2024,ma2024,bian2025,yu2015,benabou2021,qiu2025}. These reduced-order and LSTM-based approaches are generally aligned with established engineering workflows, but are typically formulated to represent specific aspects of spatio-temporal behaviour rather than the full coupled correlations among multiple outputs \parencite{zhang2022,kapusuzoglu2022,sun2023,wu2024,lin2024,xie2024,bian2025}. More recent operator-learning approaches, such as DeepONet and Fourier Neural Operators \parencite{li2023,li2024}, have also been proposed for spatio-temporal modelling; however, their implementation typically requires additional considerations in data structuring and operator representation, and is therefore not examined further here in order to retain a streamlined engineering-oriented framework.

Moreover, geotechnical applications present a unique challenge, where low-dimensional, static soil parameters shall be mapped to high-dimensional, time-dependent, and spatially distributed responses. This process introduces several key challenges: (1) effectively capturing dynamic correlations within complex output spaces, (2) mapping from static, low-dimensional inputs to evolving multi-output responses and (3) needing a practically deployable neural network architecture that can accommodate multiple output types within routine engineering workflows. A simplified, engineering-oriented architecture is therefore proposed to support integration within standard geotechnical workflows while retaining the capability for time-dependent response prediction. The first contribution lies in building upon established static-to-sequence modelling components and integrating them into a unified, encoder-free formulation for spatio-temporal multi-output prediction. The emphasis is placed on formulation and application context rather than on introducing a fundamentally new architectural class. Here, “simplified” denotes the absence of a dedicated encoder stage, as the low-dimensional static soil parameters do not require explicit feature compression. Static inputs are directly mapped into a shared latent representation, from which temporal evolution is captured using an LSTM and spatial responses are reconstructed through CNN-based decoders. By explicitly decoupling temporal modelling and spatial reconstruction, the formulation enables simultaneous prediction of multiple heterogeneous output fields within a single surrogate model while reducing architectural complexity.

To further validate the framework, it is crucial not only to demonstrate satisfactory performance in terms of error metrics but also to highlight its ability to tackle realistic geotechnical challenges. Geotechnical applications, often carried out in stages such as pit excavation, tunnelling, or pile loading, are inherently influenced by uncertainties in soil properties \parencite{gardoni2007, wang2010, asem2021,hu2022, buckley2023, gong2023, li2024, zhang2025-1, dai2025,he2025}. These uncertainties necessitate the adoption of probabilistic frameworks, within which observational data are progressively incorporated to refine predictions. Surrogate performance can be tested within a sequential progressive calibration process. As second novelty in this study, a triaxial loading problem, characterized by higher nonlinearity and input space, is selected. The surrogate is constructed and tested within a sequential Bayesian calibration scheme following our previous work \parencite{yang2024}, demonstrating its capability to provide real-time predictions while reducing uncertainties. This approach provides an alternative to \cite{yang2026}, rather than serving as a replacement for ROM, which retains their advantages owing to their physical interpretability. It aims to partially overcome the limitations of ROM, which are restricted to matrix-based data, by employing neural networks.

This paper is organised as follows. \Cref{sec: intro} outlines the background and motivation of the study. \Cref{approach} presents the selected components of the metamodel, describes the adopted neural network structure, and details the sequential calibration procedure. \Cref{sec:example1} demonstrates the metamodel using a benchmark offshore strip footing problem. \Cref{sec:example2} examines its performance in a sequential Bayesian inversion of a triaxial loading test. Finally, \Cref{sec:conclusion} summarises the key findings of the study.